\section{Architecture}\label{sec:arch}
\BMR\ is organised in three layers (Fig.~\ref{fig:arch}). The \emph{input layer} contains one reader per code and file type. The \emph{data model} is a set of plain objects shared by all analyses: a band object (eigenvalues $\varepsilon_{n\bm{k} s}$ in Hartree, fractional $\bm{k}$-points, weights $w_{\bm{k}}$, electrons per spin, the Fermi energies $\eF^{s}$ written by the code and the lattice in Bohr), a projection object (orbital weights $p^{\alpha}_{n\bm{k} s}$ with species, ion, angular momentum $l$, magnetic component $m$ and, when available, principal number $n$), a structure object (lattice vectors and fractional positions in \AA, species, forces and energies) and a density object (grid dimensions, lattice, $\rho$ in e\,\AA$^{-3}$, optional spin density). The \emph{analysis layer} operates only on these objects.

\begin{figure}[t]\centering
% Architecture diagram (TikZ)
\resizebox{\linewidth}{!}{\begin{tikzpicture}[node distance=4mm and 6mm, every node/.style={font=\small},
  box/.style={draw=bmc,rounded corners=2pt,fill=bmc!6,align=center,minimum height=8mm,inner sep=3pt},
  lay/.style={draw=gray!60,rounded corners=4pt,inner sep=5pt}]
\node[box] (c1) {CASTEP\\\scriptsize .bands .pdos\_bin .castep .cell .den\_fmt};
\node[box,right=of c1] (c2) {VASP\\\scriptsize vasprun.xml KPOINTS CHGCAR};
\node[box,right=of c2] (c3) {SIESTA\\\scriptsize .bands .EIG .PDOS .RHO .out};
\node[box,right=of c3] (c4) {QE\\\scriptsize data-file-schema.xml};
\node[box,right=of c4] (c5) {ORCA\\\scriptsize .out};
\node[box,below=9mm of c3,minimum width=13cm] (dm) {Common data model: bands (Ha), projections, structure (\AA), density (e\,\AA$^{-3}$), $E_\mathrm{F}$ as written by the code};
\node[box,below=of dm,minimum width=13cm] (an) {Analysis layer: DOS/PDOS · band edges, gaps, $m^*$ · fat bands · defects and distortion · thermodynamics · densities, Voronoi, Bader, $\Delta\rho$ · comparison · citation};
\node[box,below=9mm of an,xshift=-4.3cm] (f1) {Web application\\\scriptsize JavaScript, Plotly};
\node[box,below=9mm of an] (f2) {Desktop application\\\scriptsize same interface, offline};
\node[box,below=9mm of an,xshift=4.3cm] (f3) {Command line\\\scriptsize Python, numpy, matplotlib};
\foreach \c in {c1,c2,c3,c4,c5} \draw[-{Latex[length=2mm]},gray] (\c.south) -- (dm.north -| \c.south);
\draw[-{Latex[length=2mm]},gray] (dm) -- (an);
\foreach \f in {f1,f2,f3} \draw[-{Latex[length=2mm]},gray] (an.south -| \f.north) -- (\f.north);
\end{tikzpicture}}

\caption{Architecture of \BMR. Code-specific readers convert every output to a common data model; a single set of analyses serves three front-ends with identical semantics.}\label{fig:arch}
\end{figure}

The three front-ends are:
\begin{itemize}
  \item \textbf{Web application.} A single, self-contained HTML document with embedded JavaScript, using Plotly.js for interactive 2D/3D graphics and JSZip for archive export. Files are read locally through the File API; nothing is uploaded. Whole folders can be dropped (needed for VASP, whose files have fixed names).
  \item \textbf{Desktop application.} The same document, with all libraries and fonts bundled, displayed in a native window through pywebview \cite{pywebview} (Microsoft Edge WebView2 on Windows, Qt WebEngine or WebKitGTK on Linux, WebKit on macOS). When no native engine is available it falls back to the default browser, served on the loopback interface only. A Windows installer (Inno Setup) and stand-alone executables (PyInstaller \cite{pyinstaller}) are built by continuous integration.
  \item \textbf{Command line.} A single Python file (\code{brumare.py}, installable with \code{pip install brumare}) depending only on numpy \cite{harris2020numpy} and matplotlib \cite{hunter2007matplotlib}, so that it can be copied to a cluster and run inside batch jobs. It writes PNG/PDF/SVG figures, tab-separated \code{.dat} tables for Origin or gnuplot, VESTA-ready CHGCAR files and text reports.
\end{itemize}

\paragraph{Parity between front-ends.} The JavaScript and Python implementations are written independently but follow the same specification and are cross-checked on the same inputs (Section~\ref{sec:valid}). In Python, non-CASTEP inputs are registered in an in-memory \emph{virtual file registry}: the readers for VASP, SIESTA, QE and ORCA produce objects indexed by the file names the CASTEP pipeline expects (\code{seed.bands}, \code{seed.pdos\_bin}, \ldots), and every CASTEP reader consults the registry before the disk. This design added five codes without modifying the analysis pipeline, and allowed us to verify that the CASTEP results remained byte-identical after each extension.
